% !TeX root=../main.tex
\chapter{ طراحی و ساخت حسگر لمسی بارومتریک}
%\thispagestyle{empty} 
\section{مقدمه}
در راستای هدف اصلی این پژوهش، ابتدا یک نمونه ی اولیه حسگر لامسه مبتنی بر مکانیزم بارومتریک  طراحی و ساخته شد. این انتخاب به دلیل مزایای کلیدی متعددی، از جمله دقت بالا، قیمت مناسب، ابعاد کوچک و سادگی در پیاده‌سازی الکترونیکی، صورت گرفته است. هسته اصلی نمونه اولیه، حسگر فشار 
\LR{MEMS}\LTRfootnote{Micro-electromechanical systems}
 مدل
 \LR{MPL3115A2}
   است. این حسگرهای دیجیتال برای این کاربرد بسیار مناسب هستند، زیرا فرآیندهای حساس پردازش سیگنال، مانند تقویت و تبدیل آنالوگ به دیجیتال، را به طور مستقیم در داخل تراشه انجام می‌دهند. این پردازش داخلی، پیچیدگی طراحی سخت‌افزاری را به شدت کاهش داده و امکان یکپارچه‌سازی ساده با استفاده از پروتکل‌های ارتباطی استاندارد مانند 
\LR{I2C}\LTRfootnote{Inter-Integrated Circuit}
    را فراهم می‌آورد. بخش مهمی از روش‌شناسی ما، توسعه اولیه یک نمونه تک‌حسگری بود، که امکان‌سنجی حسگر انتخابی و ارزیابی دقیق پارامترهای لایه الاستیک، از جمله ویسکوزیته و سختی آن را فراهم کرد. این کار پایه، برای طراحی پیشرفته‌تر حیاتی بود. طراحی کلی بر اساس اصولی که در مقالات
\cite{tenzer_feel_2014,hamed_design_2020,kim_experimental_2023,kim_barotac_2022}
 بیان شده، انجام شد و عملایت های پیشرفته تر، مانند اصلاحات مکانیکی و تغییر قطر سوراخ حسگر بارومتر با الهام از 
 \cite{koiva_barometer-based_2020}
 و آزمون و خطا ادامه یافت. علاوه بر اصلاحات مکانیکی، نحوه ی اجرای لایه‌ی الاستیک هم مورد بررسی قرار گرفت و درصد و نحوه‌ی ترکیب و قالبگیری مواد اعتبار سنجی شد. نتیجه‌ی این بررسی ها منجر به انتشار یک مقاله در سی و سومین کنفرانس بین المللی برق گشت که در آن علاوه بر معرفی دستاورد های پژوهش، تحلیل ها و نمودار های مرتبط نیز به طور کامل منتشر شد.
 \cite{sheykh2025Design}
 \\
 در ادامه، پژوهش به سمت طراحی یک آرایه حسگر متشکل از سه بارومتر برای اندازه‌گیری نیروهای چند ‌محوره و مماسی گسترش یافت. این انتخاب طراحی، محدودیت اصلی نمونه‌ی اولیه‌ی تک حسگره، یعنی ناتوانی در اندازه‌گیری نیرو‌های مماسی را بر طرف می‌کرد و در عین حال مزایایی مانند دقت را نیز بهبود می‌داد.
 \cite{kim_barotac_2022}
این فصل به شرح دقیق فرآیندهای طراحی، ساخت و تحلیل عملکرد این حسگر بارومتریک می‌پردازد و یک مرور جامع از توسعه آن، از نمونه اولیه تک‌المان تا یک سیستم چندمحوری، ارائه خواهد داد.

\section{طراحی و انتخاب اجزا}
در این قسمت، دو بخش اصلی نمونه‌ی اولیه‌ی تک حسگر یعنی قطعات و ماده‌ی الاستیک معرفی و شرح داده خواهد شد.
\subsection{حسگر فشار \LR{MPL3115A2}}
\begin{figure}[t]
	\centering
	\centerline{\includegraphics[width=0.6\textwidth]{MPL3115A2}}
	\caption{حسگر فشار و ارتفاع 
		\LR{MPL3115A2}
		. }
	
\end{figure}
بخش اصلی نمونه اولیه حسگر لمسی، تراشه حسگر فشار \LR{MEMS} با مدل \LR{MPL3115A2} می‌باشد که مشخصات فنی آن در پیوست \ref{app:refMan} قرار داده شده است. این قطعه، یک حسگر فشار میکرومکانیکی است که بر پایه تغییر مقاومت المان پیزومقاومتی در یک دیافراگم سیلیکونی کار می‌کند. هنگامی که فشار بر این دیافراگم اعمال می‌شود، تغییر شکل آن باعث تغییر مقاومت الکتریکی در المان‌های پیزومقاومتی می‌شود که این تغییر به عنوان یک سیگنال الکتریکی قابل اندازه‌گیری تبدیل می‌شود.   
از منظر عملکردی، این حسگر به واسطه پاسخ پایدار خود به فشارهای استاتیک، نقش گیرنده‌های مرکل 
\LTRfootnote{Merkel’s disks}
 در پوست انسان را شبیه‌سازی می‌کند. گیرنده‌های مرکل مسئول درک نیروهای ثابت و استاتیک هستند که اطلاعات مربوط به فشار گرفتن پایدار را به مغز منتقل می‌کنند. بنابراین، استفاده از حسگر بارومتریک برای اندازه‌گیری نیروی عمودی، تقلیدی مستقیم از این عملکرد زیستی است که یکی از اصول بنیادین این پژوهش محسوب می‌شود. در جدول \ref{tab:baro_sensors} مدل های دیگری از حسگر های بارومتریک که در مقالات مشابه استفاده شده اند، بیان شده است. دلایل انتخاب حسگر \LR{MPL3115A2}
در این پژوهش شامل دقت بالا، سرعت قابل قبول و موجود بودن در بازار ایران می‌شود.
\\
\begin{table}[ht]
	\centering
	\caption{.حسگر بارومتریک استفاده شده در کار های مشابه}
	\label{tab:baro_sensors}
	\onehalfspacing
	\begin{tabular}{|r|r|r|}
		\hline
		\textbf{مرجع} & \textbf{نام حسگر}  & \textbf{پروتکل}  \\ \hline \hline
		\cite{declercq_soft_2022} & \LR{MS5840-02BA} & \LR{I2C} \\ \hline
		\cite{hamed_design_2020} & \LR{MPL115A2} & \LR{I2C} \\ \hline
		\cite{kim_artificial_2020} & \LR{MS5611} & \LR{I2C/SPI\footnotemark} \\ \hline
		\cite{kim_barotac_2022} & \LR{BMP384} & \LR{I2C/SPI} \\ \hline
		\cite{kim_experimental_2023} & \LR{MPL115A2} & \LR{I2C} \\ \hline
		\cite{koiva_barometer-based_2020} & \LR{BMP388} & \LR{I2C/SPI} \\ \hline
		\cite{tenzer_feel_2014} & \LR{MPL115A2} & \LR{I2C} \\ \hline
		\cite{shayan_shrewdshoe_2018} & \LR{MD-PS002} & \LR{Analog} \\ \hline
	\end{tabular}
\end{table}
\LTRfootnotetext[4]{Serial Peripheral Interface}
\\
\subsection{پوشش منعطف}
برای انتقال مؤثر نیرو از محیط خارجی به دیافراگم حسگر بارومتریک، و همچنین برای ایجاد یک رابط نرم و منعطف، از سیلیکون
 \LR{RTV2} \LTRfootnote{Room Temperature Vulcanizing}
 به عنوان ماده‌ی قالب‌گیری استفاده شد. این ماده به دلیل خواص کشسانی مناسب، مقرون به صرفه بودن، سهولت در استفاده، اثربخشی اثبات‌شده در پژوهش‌های پیشین و همچنین موجود بودن در بازار ایران، گزینه‌ای ایده‌آل بوده است. در فرآیند ساخت، این سیلیکون با هاردنر و یک رقیق‌کننده با نسبت‌های مشخص بر اساس جدول 
 \ref{tab:rubber_prop}
 ترکیب شد تا ویسکوزیته آن برای نفوذ به فضاهای باریک و قالب‌گیری دقیق کاهش یابد. این نسبت بر اساس آزمون و خطا از مقدار اولیه ی توصیه شده توسط سازنده به دست آمده است.

انتخاب یک پوشش نرم و کشسان، صرفاً یک تصمیم فنی نبوده، بلکه یک تصمیم الهام‌گرفته از طبیعت می‌باشد. همانطور که در فصل اول اشاره شد، نوک انگشتان انسان از بافت نرم و ارتجاعی ساخته شده که مزایای مکانیکی مهمی را فراهم می‌کند. این بافت نرم به هنگام تماس با اشیاء تغییر شکل داده و سطح تماس را افزایش می‌دهد، که این امر باعث توزیع نیرو روی ناحیه‌ای بزرگ‌تر شده و خطر آسیب به اشیاء شکننده را کاهش می‌دهد. علاوه بر این، نرمی پوشش به حسگر کمک می‌کند تا با ناهمواری‌ها و اشکال نامنظم اشیاء سازگار شود و نیاز به الگوریتم‌های پیچیده کنترلی برای جایگیری دقیق را کاهش می‌دهد. بنابراین، استفاده از سیلیکون با سختی شور\LTRfootnote{Shore}
  \LR{25A}
   در این پژوهش، تلاشی آگاهانه برای شبیه‌سازی این ویژگی سودمند فیزیکی دست انسان است.   

\begin{table}[ht]
	\centering
	\caption{.جنس لایه‌ی کشسان استفاده شده در کار های مشابه}
	\label{tab:baro_sensors}
	\onehalfspacing
	\begin{tabular}{|r|r|r|}
		\hline
		\textbf{مرجع} & \textbf{مدل}  & \textbf{سختی شور}  \\ \hline \hline
		\cite{declercq_soft_2022} & - & \LR{15A} \\ \hline
		\cite{hamed_design_2020} & \LR{RTV2} & \LR{25A} \\ \hline
		\cite{kim_artificial_2020} & \LR{silicone} & - \\ \hline
		\cite{kim_barotac_2022} & \LR{VytaFlex} & \LR{20A} \\ \hline
		\cite{kim_experimental_2023} & \LR{VytaFlex} & \LR{20A} \\ \hline
		\cite{koiva_barometer-based_2020} & \LR{SF13(RTV2)} & \LR{13ShA} \\ \hline
		\cite{tenzer_feel_2014} & \LR{VytaFlex} & \LR{20A} \\ \hline
		\cite{shayan_shrewdshoe_2018} & \LR{RTV2} & \LR{25A} \\ \hline
	\end{tabular}
\end{table}
\section{فرایند ساخت نمونه اولیه}
فرآیند ساخت حسگر لمسی بارومتریک، شامل یک روند دقیق و چندمرحله‌ای است که هدف آن ایجاد یک ساختار یکپارچه و ضد هوا برای اطمینان از اندازه‌گیری دقیق نیرو است. این فرآیند در    

شکل ۱ به تصویر کشیده شده است.
\begin{figure}[t]
	\centering
	\centerline{\includegraphics[width=0.9\textwidth]{GraphicalAbstract}}
	\caption{نمای کلی از فرآیند ساخت و کالیبراسیون حسگر:
		 (الف) برد 
		\LR{MPL3115A2}
		 و سیلیکون 
		 \LR{RTV2}
		 ، (ب) فرآیند قالب‌گیری اولیه که در آن سیلیکون کمی بالاتر از برد ریخته می‌شود؛ بلافاصله پس از ریختن سیلیکون، مجموعه برای حذف حباب‌های هوا به داخل محفظه خلاء منتقل می‌شود، (ج) محفظه فلزی به دقت حرارت داده می‌شود تا چسب آن شل شده و برداشته ‌شود. سپس سوراخ روی محفظه فلزی با مته ۳ میلی‌متر بزرگ می‌شود، (د) فرآیند قالب‌گیری سیلیکون برای ایجاد کامل بخش کشسان تکرار می‌شود و (ه) راه‌اندازی کالیبراسیون برای اعمال نیرو و جمع‌آوری داده.}
	}
\Label{fig:GA}
\end{figure}
فرآیند ساخت حسگر لمسی بارومتریک، شامل یک روند دقیق و چندمرحله‌ای است که هدف آن ایجاد یک ساختار یکپارچه و ضد هوا برای اطمینان از اندازه‌گیری دقیق نیرو است. این فرآیند از قرار دادن برد حسگر MPL3115A2 درون یک قالب مخصوص که به طور سفارشی طراحی شده، آغاز می‌شود. تلاش‌های اولیه با سیلیکون RTV2 به دلیل چسبندگی ضعیف بین سیلیکون و برد PCB، منجر به نشت هوا در فشارهای اندک می‌شد؛ برای حل این مشکل، تصمیم گرفته شد که کل ماژول حسگر، و نه فقط برد PCB، به طور کامل در سیلیکون جاسازی شود. برای اطمینان از جریان مناسب سیلیکون و پر شدن فضاهای باریک، ابتدا سیلیکون RTV2 با هاردنر و یک رقیق‌کننده به نسبت‌های مشخص (   

۹۵.۲٪
سیلیکون، 
۳.۳٪
هاردنر و 
۱.۹٪
رقیق‌کننده) ترکیب می‌شود تا ویسکوزیته آن به طور چشمگیری کاهش یابد. در مرحله قالب‌گیری اولیه، مخلوط به دقت روی کیس فلزی حسگر ریخته می‌شود تا فضای داخلی آن را پر کند. بلافاصله پس از ریختن، قالب در یک محفظه خلاء قرار می‌گیرد تا فرآیند حیاتی «دی‌گسینگ» (حذف گاز) انجام شود. این مرحله برای حذف کامل حباب‌های هوای محبوس بین سیلیکون و دیافراگم حسگر ضروری است؛ زیرا وجود هوا می‌تواند به عنوان یک عایق مکانیکی عمل کرده و انتقال فشار را کاهش دهد که منجر به کاهش دقت و نشت هوا تحت فشار می‌شود. پس از اتمام دی‌گسینگ، سیلیکون به مدت ۴۸ ساعت برای سفت شدن کامل رها می‌شود. در ادامه، یک نوآوری کلیدی این پژوهش اجرا می‌شود: کیس فلزی حسگر به آرامی با یک هویه با دمای پایین حرارت داده می‌شود تا چسب آن نرم شده و بتوان آن را با دقت از برد جدا کرد. سپس، سوراخ (orifice) موجود روی کیس فلزی به آرامی با استفاده از یک مته ۲ میلی‌متر بزرگ می‌شود. این دستکاری مکانیکی، فضای خالی بین لایه سیلیکونی و دیافراگم داخلی حسگر را به حداقل می‌رساند و کوپلینگ مکانیکی را به شکل چشمگیری بهبود می‌بخشد. این روش، در مقایسه با رویکردهای پیشین برای برداشتن محافظ حسگر، ایمن‌تر و قابل‌اعتمادتر است. با بزرگ شدن سوراخ، فشار وارد شده از سطح بیرونی به شکل مؤثرتر و با دقت بالاتری به سیگنال الکتریکی تبدیل می‌شود، که نه تنها حساسیت حسگر را به شدت افزایش می‌دهد، بلکه خطی بودن پاسخ آن و مقاومت در برابر نویزهای فیزیکی را نیز بهبود می‌بخشد. در نهایت، فرآیند قالب‌گیری سیلیکون تکرار می‌شود و این بار سیلیکون به طور کامل بر روی حسگر و تا ارتفاع ۸ میلی‌متر ریخته می‌شود. پس از اتمام این مرحله، کل مجموعه به مدت ۴۸ ساعت در دمای اتاق رها می‌شود تا سیلیکون به طور کامل سفت شده و از استحکام ساختاری آن اطمینان حاصل شود.

\section{مروری بر ادبیات موضوع}
در این قسمت باید به کارهای مشابه دیگران در گذشته اشاره کرد و وزن بیشتر این قسمت بهتر است به مقالات ژورنالی سال‌های اخیر (۲ تا ۳ سال) تخصیص داده شود. به نتایج کارهای دیگران با ذکر دقیق مراجع باید اشاره شده و جایگاه و تفاوت تحقیق شما نیز با کارهای دیگران مشخص شود. استفاده از مقالات ژورنال‌های معتبر در دو یا سه سال اخیر، می‌تواند به اعتبار کار شما بیافزاید.

\section{نتیجه‌گیری}
‌در نتیجه‌گیری آخر این فصل، با توجه به بررسی انجام شده بر روی مراجع تحقیق، بخش‌های قابل گسترش و تحقیق در آن حیطه و چشم‌اندازهای تحقیق مورد بررسی قرار می‌گیرند.	در برخی از تحقیقات، نتیجه نهایی فصل روش تحقیق، ارائهٔ یک چارچوب کار تحقیقی 
\lr{(research framework)}
است.